% Book 3, Year 1, chapter 19 -- sorting.
% Provenance (never printed): sources/book3/BRIEFS.md, chapter 19.
% ledger: b1:libstdcxx-introsort, b1:cpython-listsort, b1:cpp-sort-complexity
\chapter{Sorting}\label{ch:b1:sorting}

Before electronic computers, a population census was sorted by machines that read punched cards. A card carried one
person's answers as holes in columns; the sorting machine read one column and dropped each card into one of ten pockets,
by the digit punched there. To sort by a three-digit code, the operator ran the deck through the machine three times,
\emph{starting with the last digit}, stacking the pockets in order after each pass. It looks backwards, and it works:
this chapter proves why. It also builds the \omterm{def:b1:sorting:comparison-sort}{comparison sorts} used today, \omterm{def:b1:sorting:quicksort}{quicksort} and \omterm{def:b1:sorting:merge-sort}{merge sort}, counts their
comparisons exactly or on average, and proves that no sort that only compares keys can do much better.

\begin{recall}
Lomuto's partition rearranges a slice around a \omterm{def:b1:sorting:quicksort}{pivot}, with a proved invariant
(\cref{prop:b1:program-correctness:partition}).
Recurrences such as $T(n) = 2T(n/2) + n$ and the harmonic numbers $H_n = 1 + \frac12 + \dots + \frac1n \approx \ln n$
were met in \cref{ch:b1:asymptotic-analysis}. \omterm{met:b1:heaps:heapsort}{Heapsort} sorts \omterm{def:b1:sorting:comparison-sort}{in place} in $\Theta(n \log n)$ (\cref{met:b1:heaps:heapsort}).
Insertion sort inserts each element into the sorted prefix before it.
\end{recall}

\section{What a sort promises}\label{sec:b1:sorting:promises}

\begin{definition}[Comparison sort, in-place sort]\label{def:b1:sorting:comparison-sort}
A \emph{comparison sort}\index{comparison sort} learns about the keys only by comparing two of them. A sort is
\emph{in place}\index{in-place sort} when it uses $O(1)$ memory beyond the \omterm{def:b1:pointers-and-arrays:array}{array}, or $O(\log n)$ for a recursion stack.
\end{definition}

\begin{definition}[Stable sort]\label{def:b1:sorting:stable}
A sort is \emph{stable}\index{stable sort} when elements with equal keys keep, in the output, the order they had in the
input.
\end{definition}

Stability matters for records sorted by one field after another: sorting books by title, then stably by genre, gives
each genre's books in alphabetical order.

\begin{definition}[Inversion]\label{def:b1:sorting:inversion}
An \emph{inversion}\index{inversion (of a sequence)} of a sequence $a$ is a pair of positions $i < j$ with $a_i > a_j$.
\end{definition}

A sorted sequence has no \omterm{def:b1:sorting:inversion}{inversion}; a reversed one has all $n(n - 1)/2$. Each exchange of two neighbours removes exactly
one \omterm{def:b1:sorting:inversion}{inversion}, which is why insertion sort, which only exchanges neighbours, costs as many exchanges as the input has
\omterm{def:b1:sorting:inversion}{inversions} (\cref{exo:b1:sorting:8}).

\omimg{book3/ai/b1-sorting-returned-books.jpg}{0.45\linewidth}{Returned books go back one at a time into shelves
that are already in order: insertion sort, which is fast when only a few items are out of place.}

\section{Quicksort}\label{sec:b1:sorting:quicksort}

\begin{definition}[Quicksort, pivot, three-way partition]\label{def:b1:sorting:quicksort}
\emph{Quicksort}\index{quicksort} chooses an element, the \emph{pivot}\index{pivot}, partitions the \omterm{def:b1:pointers-and-arrays:array}{array} into the
elements smaller than the pivot, the pivot, and the others, and sorts the two parts recursively. A \emph{three-way
partition}\index{three-way partition (of an array)} separates the elements smaller than, equal to and larger than the
pivot.
\end{definition}

\omcode{code/bachelor-1/19-sorting/listings/sorts.c}{16}{29}{Lomuto's partition and the basic
\omterm{def:b1:sorting:quicksort}{quicksort}}\label{lst:b1:sorting:quicksort}

\begin{proposition}[Quicksort is correct]\label{prop:b1:sorting:qs-correct}
\code{qs_basic(a, lo, hi)} sorts $a[lo..hi)$ and terminates.
\end{proposition}

\begin{proof}
By strong induction on $m = hi - lo$. A slice of fewer than 2 elements is sorted. Otherwise the partition returns $k$
with $a[lo..k)$ smaller than $a[k]$ and $a(k..hi)$ at least $a[k]$ (\cref{prop:b1:program-correctness:partition}). Both parts have
fewer than $m$ elements, so the recursive calls sort them, and they only permute each part: the whole slice is then
sorted.
\end{proof}

\begin{omfigure}
\begin{tikzpicture}[a/.style={draw, rounded corners=2pt, font=\small\ttfamily, inner sep=3pt},
  p/.style={a, fill=omThm!25}, level distance=11mm, level 1/.style={sibling distance=50mm},
  level 2/.style={sibling distance=22mm}, edge from parent/.style={draw, ->, >={Stealth[length=1.6mm]}}]
\node[a] {7 2 9 4 1 8 \textbf{5}}
  child {node[a] {2 4 \textbf{1}} child {node[a] {(empty)}} child {node[a] {4 \textbf{2}}}}
  child {node[a] {9 8 \textbf{7}} child {node[a] {(empty)}} child {node[a] {8 \textbf{9}}}};
\node[p, font=\small\ttfamily] at (0,-1.1) {5};
\node[font=\scriptsize, align=left, anchor=west] at (4.2,0) {pivot in bold: the last element;\\ 6 comparisons};
\node[font=\scriptsize, anchor=west] at (5.2,-1.1) {2 comparisons each};
\end{tikzpicture}
\par\smallskip{\small \omterm{def:b1:sorting:quicksort}{Quicksort} of $[7, 2, 9, 4, 1, 8, 5]$: the first partition puts 5 (shaded) in its final place, with
2, 4, 1 on its left and 9, 8, 7 on its right; the parts are sorted in the same way. In all $6 + 2 + 1 + 2 + 1 = 12$
comparisons.}
\end{omfigure}

The partition of $m$ elements makes $m - 1$ comparisons. When the \omterm{def:b1:sorting:quicksort}{pivot} is always the largest element, as on a sorted
\omterm{def:b1:pointers-and-arrays:array}{array} with the last element as \omterm{def:b1:sorting:quicksort}{pivot}, one part is empty, and the total is $(n - 1) + (n - 2) + \dots + 1 = n(n - 1)/2$:
$\Theta(n^2)$. When the \omterm{def:b1:sorting:quicksort}{pivot} always splits evenly, the recurrence $C(n) = 2C(n/2) + n - 1$ gives $\Theta(n \log n)$.
What happens on average?

\begin{theorem}[Average comparisons of quicksort]\label{thm:b1:sorting:qs-average}
Over the $n!$ orders of $n$ distinct keys, all equally likely, \omterm{def:b1:sorting:quicksort}{quicksort} makes on average $C(n) = 2(n + 1)H_n - 4n$
comparisons, about $2n \ln n \approx 1.39\, n \log_2 n$.
\end{theorem}

\begin{proof}[Proof of the recurrence]
The \omterm{def:b1:sorting:quicksort}{pivot} is the last element; in exactly $(n - 1)!$ of the $n!$ orders it has rank $k + 1$, for each $k = 0, \dots, n -
1$, leaving $k$ smaller and $n - 1 - k$ larger elements. Lomuto's partition keeps the relative order of the smaller
elements, and of the larger ones up to a fixed rearrangement, so over these $(n - 1)!$ orders each part comes in all its
orders equally often. Hence $C(0) = C(1) = 0$ and
\[ C(n) = n - 1 + \frac1n \sum_{k=0}^{n-1} \bigl(C(k) + C(n - 1 - k)\bigr) = n - 1 + \frac2n \sum_{k=0}^{n-1} C(k). \]
\Cref{exo:b1:sorting:10} solves it.
\end{proof}

\omoutput{code/bachelor-1/19-sorting/outputs/qs_average.txt}{The exact average over all orders, by enumeration, against
the formula, in exact fractions.}\label{out:b1:sorting:qs-average}

The enumeration only checks the claim for $n \leq 8$, and only for this program; the proof covers every $n$. The
fair-order assumption is the weak point: input is often sorted, or nearly. Three changes make \omterm{def:b1:sorting:quicksort}{quicksort} robust. The
\omterm{def:b1:sorting:quicksort}{pivot} is the median of the first, middle and last elements, which is good on sorted input. The recursion goes into the
smaller part and a loop handles the larger one, so the stack holds at most $\log_2 n$ frames, since each frame's part is
at most half the previous one. Parts of 16 elements or fewer are left to insertion sort, which is faster on small
\omterm{def:b1:pointers-and-arrays:array}{arrays}.

\omcode{code/bachelor-1/19-sorting/listings/sorts.c}{33}{44}{\omterm{def:b1:sorting:quicksort}{Quicksort} with median of three, a short stack and insertion
sort}\label{lst:b1:sorting:quicksort-good}

Equal keys remain a trap: with all keys equal, Lomuto's partition puts every element on one side
(\cref{exo:b1:sorting:5}). The \omterm{def:b1:sorting:quicksort}{three-way partition} puts them in the middle, which is left out of the recursion.

\omcode{code/bachelor-1/19-sorting/listings/sorts.c}{48}{57}{Three-way \omterm{def:b1:sorting:quicksort}{quicksort}}\label{lst:b1:sorting:quicksort3}

\section{Merge sort}\label{sec:b1:sorting:merge}

\begin{definition}[Merge, merge sort]\label{def:b1:sorting:merge-sort}
To \emph{merge}\index{merge (of sorted sequences)} two sorted sequences is to build the sorted sequence of their
elements by repeatedly taking the smaller of their two first elements. \emph{Merge sort}\index{merge sort} sorts the two
halves of an \omterm{def:b1:pointers-and-arrays:array}{array} recursively and merges them.
\end{definition}

\omcodes{code/bachelor-1/19-sorting/listings/sorts.c}{61-67,84-91}{The \omterm{def:b1:sorting:merge-sort}{merge}, and bottom-up \omterm{def:b1:sorting:merge-sort}{merge sort}}\label{lst:b1:sorting:merge}

A \omterm{def:b1:sorting:merge-sort}{merge} of $m$ elements makes at most $m - 1$ comparisons. The top-down version splits recursively; the bottom-up
version \omterm{def:b1:sorting:merge-sort}{merges} runs of width 1, then 2, 4, and so on, with no recursion. Both make at most $W(n) = W(\lfloor n/2 \rfloor)
+ W(\lceil n/2 \rceil) + n - 1$ comparisons, which is $n \log_2 n - n + 1$ for $n = 2^k$ (induction: $2(2^{k-1}(k - 1) -
2^{k-1} + 1) + 2^k - 1 = k2^k - 2^k + 1$). On a tie the \omterm{def:b1:sorting:merge-sort}{merge} takes the element of the left run first, which makes
\omterm{def:b1:sorting:merge-sort}{merge sort} \omterm{def:b1:sorting:stable}{stable}. Its price is an extra \omterm{def:b1:pointers-and-arrays:array}{array} of $n$ elements.

\omoutput{code/bachelor-1/19-sorting/outputs/stable.txt}{The library's books sorted by title, then by genre: the poetry
shelf.}\label{out:b1:sorting:stable}

\section{The lower bound}\label{sec:b1:sorting:lower-bound}

\begin{definition}[Comparison tree]\label{def:b1:sorting:comparison-tree}
The \emph{comparison tree}\index{comparison tree} of a \omterm{def:b1:sorting:comparison-sort}{comparison sort} for inputs of size $n$ is the \omterm{def:b1:induction-and-termination:binary-tree}{binary tree} whose
\omterm{def:b1:induction-and-termination:rooted-tree}{internal nodes} are the comparisons the sort makes, with one \omterm{def:b1:induction-and-termination:rooted-tree}{subtree} for each outcome, and whose leaves are the outputs.
An input follows one path from the root to a \omterm{def:b1:induction-and-termination:rooted-tree}{leaf}; the length of the path is the number of comparisons.
\end{definition}

\begin{omfigure}
\begin{tikzpicture}[q/.style={draw, rounded corners=2pt, font=\small, inner sep=2pt, fill=omCode!8},
  l/.style={font=\small, inner sep=2pt}, level distance=11mm, level 1/.style={sibling distance=56mm},
  level 2/.style={sibling distance=28mm}, level 3/.style={sibling distance=16mm},
  e/.style={font=\scriptsize, midway, fill=white, inner sep=1pt}]
\node[q] {$b < a$?}
  child {node[q] {$c < a$?}
    child {node[q] {$c < b$?} child {node[l] {$c\,b\,a$} edge from parent node[e] {yes}}
      child {node[l] {$b\,c\,a$} edge from parent node[e] {no}} edge from parent node[e] {yes}}
    child {node[l] {$b\,a\,c$} edge from parent node[e] {no}} edge from parent node[e] {yes}}
  child {node[q] {$c < b$?}
    child {node[q] {$c < a$?} child {node[l] {$c\,a\,b$} edge from parent node[e] {yes}}
      child {node[l] {$a\,c\,b$} edge from parent node[e] {no}} edge from parent node[e] {yes}}
    child {node[l] {$a\,b\,c$} edge from parent node[e] {no}} edge from parent node[e] {no}};
\end{tikzpicture}
\par\smallskip{\small The \omterm{def:b1:sorting:comparison-tree}{comparison tree} of insertion sort on three keys $a, b, c$: six leaves, one for each order, and
height 3. A sorted input takes the rightmost path, two comparisons.}
\end{omfigure}

\begin{theorem}[Lower bound for comparison sorts]\label{thm:b1:sorting:lower-bound}
Every \omterm{def:b1:sorting:comparison-sort}{comparison sort} makes at least $\lceil \log_2 n! \rceil$ comparisons on some input of $n$ distinct keys, and
$\log_2 n! \geq \frac n2 \log_2 \frac n2$: the worst case is in $\Omega(n \log n)$.
\end{theorem}

\begin{proof}
The $n!$ orders of $n$ distinct keys need $n!$ different outputs, so the \omterm{def:b1:sorting:comparison-tree}{comparison tree} has at least $n!$ leaves. A
\omterm{def:b1:induction-and-termination:binary-tree}{binary tree} of height $h$ has at most $2^h$ leaves, so $2^h \geq n!$ and $h \geq \lceil \log_2 n! \rceil$. Finally, $n!$
has at least $n/2$ factors at least $n/2$, so $n! \geq (n/2)^{n/2}$.
\end{proof}

\omoutput{code/bachelor-1/19-sorting/outputs/lower_bound.txt}{The bound $\lceil \log_2 n! \rceil$ against the worst cases
of three sorts.}\label{out:b1:sorting:lower-bound}

\omterm{def:b1:sorting:merge-sort}{Merge sort} stays within about 10\,\% of the bound; \omterm{def:b1:sorting:merge-sort}{merge} insertion, the method of the problem, meets it exactly up to $n =
11$, then is off by one. The same argument bounds the average: a \omterm{def:b1:induction-and-termination:binary-tree}{binary tree} with $N$ leaves has an average \omterm{def:b1:induction-and-termination:rooted-tree}{leaf} depth at
least $\log_2 N$, so even on average a \omterm{def:b1:sorting:comparison-sort}{comparison sort} makes at least $\log_2 n!$ comparisons. \omterm{def:b1:sorting:quicksort}{Quicksort}'s $1.39\, n
\log_2 n$ is, for large $n$, 39\,\% above it.

\begin{omfigure}
\begin{tikzpicture}
\begin{axis}[width=11cm, height=6.2cm, xmin=0, xmax=66, ymin=0, ymax=600, xlabel={number of keys $n$},
  ylabel={comparisons}, legend pos=north west, legend style={font=\small}, tick label style={font=\small}]
\addplot[thick, omThm, dashed] table[x=n, y=qssorted] {code/bachelor-1/19-sorting/figures/comparisons.dat};
\addlegendentry{quicksort, sorted input}
\addplot[thick, omThm, mark=*, mark size=1pt] table[x=n, y=qsrandom]
  {code/bachelor-1/19-sorting/figures/comparisons.dat}; \addlegendentry{quicksort, random}
\addplot[thick, omDef, mark=square*, mark size=1pt] table[x=n, y=merge]
  {code/bachelor-1/19-sorting/figures/comparisons.dat}; \addlegendentry{merge sort, random}
\addplot[thick, omCode, mark=triangle*, mark size=1.2pt] table[x=n, y=heap]
  {code/bachelor-1/19-sorting/figures/comparisons.dat}; \addlegendentry{heapsort, random}
\addplot[thick, black, densely dotted] table[x=n, y=bound] {code/bachelor-1/19-sorting/figures/comparisons.dat};
\addlegendentry{$\lceil \log_2 n! \rceil$}
\end{axis}
\end{tikzpicture}
\par\smallskip{\small Comparisons from \cref{out:b1:sorting:counts}, averaged over 200 random orders. \omterm{def:b1:sorting:quicksort}{Quicksort} on sorted
input leaves the chart at $n = 36$ (630 comparisons). \omterm{def:b1:sorting:merge-sort}{Merge sort} hugs the bound; heapsort makes about twice as many.}
\end{omfigure}

\omoutput{code/bachelor-1/19-sorting/outputs/counts.txt}{Comparisons of the sorts of \cref{lst:b1:sorting:quicksort,lst:b1:sorting:quicksort-good,lst:b1:sorting:merge}
and of heapsort; \code{qs+} is \cref{lst:b1:sorting:quicksort-good}.}\label{out:b1:sorting:counts}

\section{Sorting without comparisons}\label{sec:b1:sorting:radix}

The lower bound holds only for algorithms that compare keys. Keys that are small integers, or strings of digits, can be
sorted by looking at their digits.

\begin{definition}[Counting sort, radix sort]\label{def:b1:sorting:radix}
\emph{Counting sort}\index{counting sort} sorts $n$ keys in $\{0, \dots, k - 1\}$ by counting the keys of each value,
computing from the counts where each value's block starts, and copying each key to the next free cell of its block, in
input order. \emph{Radix sort}\index{radix sort}, least significant digit first, sorts keys of $d$ digits by $d$
passes of a \omterm{def:b1:sorting:stable}{stable sort} on one digit, from the last digit to the first.
\end{definition}

\omcode{code/bachelor-1/19-sorting/listings/radix.c}{6}{12}{\omterm{def:b1:sorting:radix}{Counting sort} on one decimal digit}\label{lst:b1:sorting:counting}

\omterm{def:b1:sorting:radix}{Counting sort} costs $\Theta(n + k)$ and is \omterm{def:b1:sorting:stable}{stable}, since keys are copied in input order into their block. \omterm{def:b1:sorting:radix}{Radix sort}
costs $\Theta(d(n + k))$: linear in $n$ for keys of a fixed number of digits.

\omoutput{code/bachelor-1/19-sorting/outputs/radix_demo.txt}{\omterm{def:b1:sorting:radix}{Radix sort} of seven 3-digit keys, one pass per
digit.}\label{out:b1:sorting:radix-demo}

\begin{omfigure}
\begin{tikzpicture}[k/.style={font=\small\ttfamily, inner sep=1.5pt}, >={Stealth[length=1.4mm]}]
\foreach \col/\keys/\hl in {0/{329,457,657,839,436,720,355}/-1, 1/{720,355,436,457,657,329,839}/2,
  2/{720,329,436,839,355,457,657}/1, 3/{329,355,436,457,657,720,839}/0} {
  \foreach \kk [count=\r from 0] in \keys {
    \node[k] (n\col-\r) at (\col*2.6, -\r*0.5) {\kk};}}
\foreach \col/\t in {0/input, 1/by units, 2/by tens, 3/by hundreds} {\node[font=\scriptsize] at (\col*2.6,0.55) {\t};}
\foreach \a/\b in {1/4, 3/5} {\draw[->, omThm, thick] (n1-\a.east) -- (n2-\b.west);}
\node[font=\scriptsize, text=omThm, align=left, anchor=north] at (3.9,-3.4) {355 and 457 both have a 5 in the tens:\\
they keep the order of the previous pass};
\end{tikzpicture}
\par\smallskip{\small The passes of \cref{out:b1:sorting:radix-demo}. The arrows follow 457 and 355 into the pass on the
tens: equal digits, so the \omterm{def:b1:sorting:stable}{stable} pass leaves them in the order the units gave them.}
\end{omfigure}

\begin{proposition}[Radix sort is correct]\label{prop:b1:sorting:radix}
After the pass on digit $j$ (counting from 1 at the units), the keys are sorted by their last $j$ digits.
\end{proposition}

The proof, by induction on $j$, rests on the stability of each pass (\cref{exo:b1:sorting:12}); with an unstable pass,
it fails (\cref{exo:b1:sorting:6}). The punched-card machines used exactly this method, one column per pass. The
returns of the town's loans, as eight-digit numbers \code{yyyymmdd}, sort in eight passes:

\omoutput{code/bachelor-1/19-sorting/outputs/radix_loans.txt}{\omterm{def:b1:sorting:radix}{Radix sort} of the loans' return
dates.}\label{out:b1:sorting:radix-loans}

\begin{history}[Card sorters]
% ledger: b1:card-sorter
The machines of this chapter's opening were built for decades. IBM's type 82 sorter of 1949 read one column of each
card with a metal brush and dropped the card into one of 13 pockets, one for each of the 12 rows a hole could be
punched in and one for blank or rejected cards, at 650 cards a minute. A deck was sorted on a key of several columns
one column at a time, least significant first, the operator collecting the pockets in order between passes: a \omterm{def:b1:sorting:radix}{radix
sort}, run by hand.
\end{history}

\omimg{book3/ibm-82-card-sorter.jpg}{0.55\linewidth}{An IBM 82 card sorter at the Computer History Museum. The row of
pockets along the front received the cards, one pocket per row of the card. Photograph:
leighklotz, CC BY 2.0.}

\section{How library sorts are built}\label{sec:b1:sorting:library}

\begin{definition}[Introsort]\label{def:b1:sorting:introsort}
\emph{Introsort}\index{introsort} is a \omterm{def:b1:sorting:quicksort}{quicksort} that counts its recursion depth and switches to heapsort on a part once
the depth passes a limit proportional to $\log n$; small parts are left to insertion sort.
\end{definition}

\omterm{def:b1:sorting:introsort}{Introsort} keeps \omterm{def:b1:sorting:quicksort}{quicksort}'s speed on ordinary input and heapsort's $O(n \log n)$ guarantee on inputs that defeat the
\omterm{def:b1:sorting:quicksort}{pivot} choice. The C++ standard requires of \code{std::sort} $O(n \log n)$ comparisons in the worst case, and allows it to
be unstable; \code{std::stable_sort} is the \omterm{def:b1:sorting:stable}{stable} one.

\begin{dated}{2026}
The GNU C++ library implements \code{std::sort} as an \omterm{def:b1:sorting:introsort}{introsort}: a depth limit of $2\lfloor \log_2 n \rfloor$, a
heap-based sort when it runs out, parts of up to 16 elements left for a final insertion sort. CPython sorts lists with
\emph{timsort}, a \omterm{def:b1:sorting:stable}{stable}, adaptive \omterm{def:b1:sorting:merge-sort}{merge sort} that finds the runs already sorted in the input and \omterm{def:b1:sorting:merge-sort}{merges} them; its \omterm{def:b1:sorting:merge-sort}{merge}
order now follows the \emph{powersort} strategy. Both are read from the libraries' source code; such details change
between versions.
\end{dated}

\omoutput{code/bachelor-1/19-sorting/outputs/timing.txt}{Sorting a million random integers, measured on the author's
laptop.}\label{out:b1:sorting:timing}

The C library's \code{qsort} calls a comparison function through a \omterm{def:b1:pointers-and-arrays:pointer}{pointer} for every comparison, which the sorts of this
chapter avoid; \omterm{def:b1:sorting:radix}{radix sort}, with nine passes over the data and no comparison, is the fastest here. Measured times depend
on the machine and the \omterm{def:b1:compilation-pipeline:compiler}{compiler}, and say nothing about the worst case.

\section{Exercises}\label{sec:b1:sorting:exercises}

\begin{exercise}[$\star$]\label{exo:b1:sorting:1}
Apply Lomuto's partition, with the last element as \omterm{def:b1:sorting:quicksort}{pivot}, to $[3, 8, 2, 5, 1, 4, 7, 6]$. Give the \omterm{def:b1:pointers-and-arrays:array}{array} and the
\omterm{def:b1:sorting:quicksort}{pivot}'s index.
\end{exercise}

\begin{exercise}[$\star$]\label{exo:b1:sorting:2}
\omterm{def:b1:sorting:merge-sort}{Merge} $[1, 4, 7, 9]$ and $[2, 3, 8, 10, 11]$, and count the comparisons.
\end{exercise}

\begin{exercise}[$\star$]\label{exo:b1:sorting:3}
\omterm{def:b1:sorting:radix}{Counting sort} is applied to $[3, 1, 4, 1, 5, 2, 0, 3]$ with $k = 6$. Give the count \omterm{def:b1:pointers-and-arrays:array}{array}, the block starts and the
output.
\end{exercise}

\begin{exercise}[$\star$]\label{exo:b1:sorting:4}
The records $(2, \mathrm{a}), (1, \mathrm{b}), (2, \mathrm{c}), (1, \mathrm{d})$ are sorted by their number. Which
outputs are those of a \omterm{def:b1:sorting:stable}{stable sort}? Is selection sort, which exchanges the minimum of the rest with the current element,
\omterm{def:b1:sorting:stable}{stable}?
\end{exercise}

\begin{exercise}[$\star\star$]\label{exo:b1:sorting:5}
Explain the counts of the basic \omterm{def:b1:sorting:quicksort}{quicksort} on $n$ equal keys, and why the three-way version makes $2n$ comparisons.
\omoutput[2-4]{code/bachelor-1/19-sorting/outputs/exo_equal.txt}{Equal keys.}
\end{exercise}

\begin{exercise}[$\star\star$]\label{exo:b1:sorting:6}
This \omterm{def:b1:sorting:radix}{counting sort} fills each digit's block from its end. Show that it sorts by one digit correctly, explain the output
of a \omterm{def:b1:sorting:radix}{radix sort} built on it, and repair it.
\omcode{code/bachelor-1/19-sorting/listings/exo_radix_bug.c}{5}{10}{Filling the blocks backwards}
\omoutput[2-2]{code/bachelor-1/19-sorting/outputs/exo_radix_bug.txt}{\omterm{def:b1:sorting:radix}{Radix sort} of the keys of
\cref{out:b1:sorting:radix-demo}.}
\end{exercise}

\begin{exercise}[$\star\star$]\label{exo:b1:sorting:7}
Sort the town's loans by member, and by date within each member, with a \omterm{def:b1:sorting:stable}{stable sort} applied twice. In which order must
the two sorts be applied? Test the result against a sort on the pair.
\end{exercise}

\begin{exercise}[$\star\star$]\label{exo:b1:sorting:8}
Prove that insertion sort makes exactly as many exchanges as its input has \omterm{def:b1:sorting:inversion}{inversions}.
\end{exercise}

\begin{exercise}[$\star\star$]\label{exo:b1:sorting:9}
Compute $\lceil \log_2 12! \rceil$. What does it say about sorting 12 keys? Compare with \omterm{def:b1:sorting:merge-sort}{merge sort} and \omterm{def:b1:sorting:merge-sort}{merge} insertion
in \cref{out:b1:sorting:lower-bound}.
\end{exercise}

\begin{exercise}[$\star\star\star$]\label{exo:b1:sorting:10}
Solve the recurrence of \cref{thm:b1:sorting:qs-average}: multiply by $n$, subtract the equation for $n - 1$, divide by
$n(n + 1)$, and sum.
\end{exercise}

\begin{exercise}[$\star\star\star$]\label{exo:b1:sorting:11}
Count the \omterm{def:b1:sorting:inversion}{inversions} of an \omterm{def:b1:pointers-and-arrays:array}{array} in $O(n \log n)$ by modifying \omterm{def:b1:sorting:merge-sort}{merge sort}, and test against the quadratic count.
\end{exercise}

\begin{exercise}[$\star\star\star$]\label{exo:b1:sorting:12}
Prove \cref{prop:b1:sorting:radix} by induction on $j$.
\end{exercise}

\section{Problem: Five keys in seven comparisons}\label{sec:b1:sorting:problem}

\begin{problem}[{Weekend problem --- Five keys in seven comparisons}]\label{pb:b1:sorting:1}
How few comparisons can sort five keys? A program sorts all 120 orders of five keys with insertion sort, \omterm{def:b1:sorting:merge-sort}{merge sort} and
\emph{\omterm{def:b1:sorting:merge-sort}{merge} insertion}, counting the comparisons with a key type that counts them.
\omcode{code/bachelor-1/19-sorting/listings/sort5.py}{33}{43}{\omterm{def:b1:sorting:merge-sort}{Merge} insertion for five keys}
\omoutput{code/bachelor-1/19-sorting/outputs/sort5.txt}{All 120 orders, for three methods.}\label{out:b1:sorting:sort5}

\medskip\noindent\textbf{Part I --- The lower bound.}
\begin{enumerate}
\item Compute $\lceil \log_2 120 \rceil$, and conclude that no \omterm{def:b1:sorting:comparison-sort}{comparison sort} sorts every order of five keys in six
comparisons.
\item Why is the bound about the worst case, and what does the average bound $\log_2 120$ say?
\item How many leaves at least has the \omterm{def:b1:sorting:comparison-tree}{comparison tree} of a sort of five keys, and how many leaves at most has a tree of
height 7?
\item Show that a \omterm{def:b1:sorting:comparison-tree}{comparison tree} of height 7 with 120 leaves has at most 8 leaves at depth 6 or less. (Extend each
\omterm{def:b1:induction-and-termination:rooted-tree}{leaf} at depth $d$ to the $2^{7 - d}$ positions below it at depth 7.) How many does \omterm{def:b1:sorting:merge-sort}{merge} insertion have?
\item After the first three comparisons of \omterm{def:b1:sorting:merge-sort}{merge} insertion, how many of the 120 orders are still possible? Why can at
most $2^4$ remain after three comparisons of any method that always finishes in seven?
\item For $n = 4$, compute the bound, and find a method meeting it.
\end{enumerate}

\medskip\noindent\textbf{Part II --- Two familiar methods.}
\begin{enumerate}[resume]
\item Which input makes insertion sort take 10 comparisons on five keys, and which 4?
\item Explain insertion sort's average, 7.72, from the \omterm{def:b1:sorting:inversion}{inversions}.
\item Compute \omterm{def:b1:sorting:merge-sort}{merge sort}'s worst case on five keys from the recurrence $W(n)$.
\item Find an order of five keys on which \omterm{def:b1:sorting:merge-sort}{merge sort} makes 8 comparisons.
\item Why does \omterm{def:b1:sorting:merge-sort}{merge sort} never take fewer than 5 comparisons on five keys?
\item Which of the two methods is closer to the bound?
\end{enumerate}

\medskip\noindent\textbf{Part III --- Merge insertion.}
\begin{enumerate}[resume]
\item After the first three comparisons, what is known about $a, b, c, d$?
\item Why does inserting $e$ into the chain $a < b < d$ cost at most 2 comparisons?
\item Why does $c$ need to be compared only with the elements below $d$, and why are there at most three of them?
\item Count the comparisons of the method in the worst case.
\item Trace the method on the order $(4, 1, 3, 0, 2)$ of the keys $0, \dots, 4$ assigned to $a, \dots, e$.
\item Why does inserting $e$ before $c$ give a better worst case than inserting $c$ first?
\end{enumerate}

\medskip\noindent\textbf{Part IV --- Checking every order.}
\begin{enumerate}[resume]
\item What does the program check on each of the 120 orders?
\item Why is checking all 120 orders a proof for five distinct keys, and why are distinct keys enough?
\item How many orders take only 6 comparisons? Is that compatible with the lower bound?
\item The general method sorts $n$ keys in $\sum_{k=1}^{n} \lceil \log_2 \frac{3k}{4} \rceil$ comparisons. Check this
for $n = 5$ and compare with the bound for $n = 12$.
\item Why is an exhaustive check impossible for $n = 20$?
\item State the named result: \emph{five keys can be sorted in seven comparisons, and no method does it in six}.
\end{enumerate}
\end{problem}
